% Documento già completo di preambolo.
%
% Aprendolo, TEX2PDF lo dichiara «documento completo» nella barra in alto e
% disattiva la galleria: il sorgente porta la propria veste, e sovrapporgliene
% una seconda produrrebbe solo errori confusi.
%
% Il pulsante «estrai il corpo» fa la conversione: prende ciò che sta fra
% \begin{document} e \end{document} come nuovo sorgente e sposta il resto del
% preambolo fra le aggiunte, restituendo l'uso dei template. L'operazione si
% annulla.

\documentclass[11pt,a4paper]{article}

\usepackage[T1]{fontenc}
\usepackage[utf8]{inputenc}
\usepackage[italian]{babel}
\usepackage[margin=25mm]{geometry}
\usepackage{mathpazo}

\title{Documento con preambolo proprio}
\author{Chi l'ha scritto}
\date{\today}

\begin{document}

\maketitle

\section{Perché la galleria è spenta}

Questo sorgente dichiara la propria classe, i propri font e i propri margini.
Il template non ha nulla da aggiungere: il documento si compila com'è.

\section{Che cosa fa l'estrazione}

Prende queste sezioni e le mette nell'editor come corpo; manda
\verb|\usepackage{mathpazo}| e le altre righe del preambolo nel pannello delle
aggiunte. Da lì in poi il documento torna a essere impaginabile con i quattro
template, e i metadati del modulo tornano a valere.

\end{document}
